\section {Regrouping coins (Combinatorics)}

This riddle is from Doron.

A collection of $\frac{n(n+1)}{2}$ identical coins is split into groups in some way.
Now repeat the following step: take one coin from every group and put all the coins you took into a new group.
A group that runs out of coins disappears.
Prove that eventually the groups will have the sizes $1, 2, 3, \ldots, n$.

\subsection{Solution}

First notice that the sizes $1, 2, \ldots, n$ do not change once they are reached: taking a coin from each of the $n$ groups
leaves groups of sizes $0, 1, \ldots, n-1$ and creates a new group of size $n$.

Draw the groups as rows of coins, sorted from the largest group at the top to the smallest at the bottom,
with all rows starting at the left. The coin in row $i$ and column $j$ sits at position $(i,j)$, and we say it is at level $i+j$.
Taking one coin from every group means taking the first column, and the new group becomes a new row.
Put the new row at the top for now, and push all the other rows down by one.
A coin at $(i,j)$ with $j \ge 2$ moves to $(i+1, j-1)$ and a coin of the first column at $(i,1)$ moves to $(1,i)$.
Either way, every coin stays on its level.
Then the new row may have to move down to its place to keep the rows sorted.

Let the energy of the coins be the sum of their levels. The step that moves the coins keeps the energy.
Moving the new row down to its place lowers the energy: it moves the new row to a larger row number,
but it moves the rows it passes, which are longer, to smaller row numbers. So the energy never grows,
and since it is a positive whole number it can only go down a finite number of times.
From some point on the new row is always at its place at the top, and the coins only move along their levels.

A level $t+1$ has $t$ positions, $(1,t), (2,t-1), \ldots, (t,1)$, and every step moves the coins on it one position along this list,
cyclically, with $(t,1)$ moving to $(1,t)$.
So from that point on, the coins on level $t+1$ simply rotate with a period of $t$ steps.

Suppose the groups never reach $1, 2, \ldots, n$. Then the levels $2, 3, \ldots, n+1$, which together have room for exactly
$\frac{n(n+1)}{2}$ coins, are not all full. Let level $d+1$ be the lowest one that is not full.
The coins that are not on the levels below it are more than it can hold, so there are coins above it, and because the rows start
at the left there is a coin on level $d+2$.
The empty position on level $d+1$ rotates with period $d$ and the coin on level $d+2$ rotates with period $d+1$.
Since $d$ and $d+1$ have no common divisor, by the Chinese remainder theorem there is a step at which both are in the same row $r$:
the empty position at $(r, d+1-r)$ and the coin right after it at $(r, d+2-r)$.
But a row cannot have a coin right after an empty position. This contradiction shows that the groups do reach the sizes $1, 2, \ldots, n$.
